\documentclass[12pt]{article}

\usepackage[a4paper, margin=2.5cm]{geometry}
\usepackage{amsmath}
\usepackage{amssymb}
\usepackage{amsthm}
\usepackage[colorlinks=true, linkcolor=blue, citecolor=blue, urlcolor=blue]{hyperref}
\usepackage{booktabs}
\usepackage{natbib}

\newtheorem{definition}{Definition}
\newtheorem{remark}{Remark}

\newcommand{\Mcl}{\mathcal{M}_{\mathrm{cl}}}
\newcommand{\Ccal}{\mathcal{C}}
\newcommand{\Fcal}{\mathcal{F}}
\newcommand{\Pcal}{\mathcal{P}}
\newcommand{\Tcal}{\mathcal{T}}
\newcommand{\Sadm}{\mathcal{S}_{\mathrm{adm}}}
\newcommand{\Scl}{\mathcal{S}_{\mathrm{cl}}}
\newcommand{\phig}{\varphi}

\title{A Correspondence Between Constraint-Manifold Geometry\\and Standard Model Structure}
\author{%
  Lee Smart\\[4pt]
  \textit{Institute of Vibrational Field Dynamics}\\[2pt]
  \texttt{contact@vibrationalfielddynamics.org}\\[2pt]
  \texttt{@vfd\_org}%
}
\date{April 2026}

\begin{document}

\maketitle

\noindent\textbf{Status:} Preprint (not peer-reviewed)

\begin{abstract}
The preceding papers in this programme (V--X) developed a constraint-based framework for particle mass, electroweak structure, confinement, continuous geometry, and a gravitational analogy, all grounded in the spectral properties of the 600-cell and the closure-operator formalism. However, this framework is not expressed in the language of the Standard Model: it uses constraints, admissibility, and geometric projection rather than fields, Lagrangians, and gauge coupling.

The present paper provides a systematic structural correspondence between the constraint-manifold framework and Standard Model concepts. A mapping table relates framework objects (admissible states, closed set, projection invariance, closure residual, connectivity constraint, manifold curvature) to their Standard Model counterparts (Hilbert space, physical states, gauge symmetry, mass, confinement, gravitational structure). Each physical sector --- mass, electroweak, strong, and gravity --- is addressed in turn.

This paper does not derive the Standard Model, reproduce its Lagrangian, or claim physical equivalence. It provides a translation layer that makes the constraint framework legible to physicists working within the Standard Model paradigm, and explicitly identifies where the correspondence holds, where it is only structural, and where it remains open.
\end{abstract}

%==================================================================
\section{Introduction}\label{sec:intro}
%==================================================================

Papers~V--X~\citep{SmartV, SmartVI, SmartVII, SmartVIII, SmartIX, SmartX} developed a constraint-based framework in which particle masses, electroweak structure, confinement, continuous geometry, and a gravitational analogy all emerge from the closure properties of a $\phig$-structured geometric manifold. The framework rests on explicit postulates, uses exact spectral data from the 600-cell, and maintains careful epistemic boundaries throughout.

However, the framework is not written in the language most physicists use. The Standard Model is formulated in terms of quantum fields, gauge symmetries, Lagrangian densities, and renormalisation. The closure framework uses constraint sets, metric projections, closure functionals, and graph-theoretic invariants.

The present paper bridges this gap. It does not derive the Standard Model or claim equivalence. It provides a \textit{structural correspondence}: a systematic mapping showing how each major concept in the Standard Model has a counterpart in the constraint framework, what the relationship is, and where the correspondence breaks down.

\subsection*{Scope and epistemic status}

This is a translation paper, not a derivation paper. Its purpose is to make the constraint framework interpretable to physicists, not to prove it correct. Every correspondence is explicitly classified as structural, analogical, or open. The paper introduces no new physics beyond what is already in Papers~V--X.

%==================================================================
\section{Conceptual Difference: Constraints vs Lagrangians}\label{sec:difference}
%==================================================================

The Standard Model and the constraint framework begin from fundamentally different starting points:

\begin{table}[ht]
\centering
\caption{Starting-point comparison.}
\label{tab:startingpoint}
\begin{tabular}{@{}ll@{}}
\toprule
Standard Model & Constraint framework \\
\midrule
Fields on spacetime & States in an admissible state space \\
Lagrangian density $\mathcal{L}$ & Closure functional $\Fcal(\psi)$ \\
Equations of motion from $\delta S = 0$ & Closure-stable states from $\Fcal(\psi) = 0$ \\
Dynamics first, constraints follow & Constraints first, dynamics deferred \\
Parameters fitted to experiment & Spectral data from the 600-cell (zero fitted) \\
Gauge symmetry postulated & Projection invariance arises as a structural feature of closure \\
\bottomrule
\end{tabular}
\end{table}

The Standard Model is a \textit{dynamical} framework: it specifies fields, their interactions, and their evolution equations. The constraint framework is a \textit{structural} framework: it specifies which states are admissible, which satisfy closure, and what geometric properties they carry. The relationship between the two is analogous to the relationship between kinematics and dynamics: the constraint framework specifies the allowed configurations, while the Standard Model specifies the evolution within those configurations.

\begin{remark}
The constraint framework does not currently include a Lagrangian, field equations, or a time-evolution law. These would need to be introduced for the framework to make dynamical predictions. The present paper addresses only the structural correspondence.
\end{remark}

%==================================================================
\section{Core Dictionary}\label{sec:dictionary}
%==================================================================

The following table provides the central structural correspondence between Standard Model concepts and constraint-framework objects:

\begin{table}[ht]
\centering
\caption{Structural correspondence: Standard Model $\leftrightarrow$ constraint framework.}
\label{tab:dictionary}
\begin{tabular}{@{}lll@{}}
\toprule
Standard Model & Constraint framework & Status \\
\midrule
Hilbert space of states & $\Sadm \subset V_\infty$ & Structural \\
Physical states & $\Mcl$ (constraint manifold) & Structural \\
Gauge orbit & Projection/torsional equivalence class & Structural \\
Gauge symmetry & Invariance under $\Pcal$ (projection) & Structural \\
Gauge bosons & Torsional connection-like modes & Analogical \\
Mass (Yukawa / Higgs) & $m = \kappa\, d(\psi, \Mcl)$ (closure residual) & Structural \\
Massless states & States on $\Mcl$ ($\Delta = 0$) & Structural \\
Massive states & States off $\Mcl$ ($\Delta > 0$) & Structural \\
Confinement & Connectivity constraint on $\Mcl$ & Structural \\
Colour charge & Threefold composite admissibility & Analogical \\
Spacetime manifold (GR) & $\Mcl$ (constraint manifold; \textit{not} spacetime) & Analogical \\
Spacetime metric $g_{\mu\nu}$ & Induced metric $g_{\Mcl}$ & Analogical \\
Geodesic motion & Constrained geodesic-like motion on smooth strata & Analogical \\
Curvature (gravity) & Constraint-interaction curvature & Analogical \\
Lagrangian / action & Not provided & Open \\
Renormalisation & Not addressed & Open \\
Scattering amplitudes & Not computable in current framework & Open \\
Running couplings & Not addressed & Open \\
\bottomrule
\end{tabular}
\end{table}

\textsc{Key:}
\begin{itemize}
    \item \textbf{Structural:} the correspondence is well-defined and formally established within the framework.
    \item \textbf{Analogical:} the correspondence is at the level of geometric/structural analogy, not formal derivation.
    \item \textbf{Open:} no correspondence has been established; this is a gap in the framework.
\end{itemize}

%==================================================================
\section{Mass Sector}\label{sec:mass}
%==================================================================

\textbf{Standard Model:} Particle masses are free parameters, set by Yukawa couplings to the Higgs field. The Higgs mechanism generates mass through spontaneous symmetry breaking. The 13+ mass parameters are fitted to experiment.

\textbf{Constraint framework (Paper~V~\citep{SmartV}):} Mass is not an input parameter but a structural quantity: the closure residual $\Delta(\psi) = d(\psi, \Mcl)$. The mass of each particle is determined by a spectral-geometric ansatz combining the closure invariant, standing-wave self-consistency corrections, and Hopf-fiber winding energies, all computed from the 600-cell graph. Under explicit postulates and structural assignments, the framework yields high-accuracy numerical correspondences for 13 non-reference particle masses with zero fitted continuous parameters.

\textbf{Correspondence status:} Structural. The framework provides an alternative structural account of mass (closure residual vs.\ Higgs field) that is compatible with observed values but does not derive the Higgs mechanism or replace it.

%==================================================================
\section{Electroweak Sector}\label{sec:ew}
%==================================================================

\textbf{Standard Model:} The electroweak interaction is governed by the gauge group $\mathrm{SU}(2)_L \times \mathrm{U}(1)_Y$. The $W$ and $Z$ bosons acquire mass through the Higgs mechanism; the photon remains massless. The Weinberg mixing angle $\theta_W$ relates the gauge eigenstates to the mass eigenstates.

\textbf{Constraint framework (Paper~VI~\citep{SmartVI}):} The electroweak gauge structure is reinterpreted as torsional invariance under projection: $\mathrm{SU}(2)_L$ corresponds to internal torsional doublet symmetry, $\mathrm{U}(1)_Y$ to phase-preserving global rotation, and gauge invariance to projection invariance ($\Pcal(\Tcal_\theta\psi) = \Pcal(\psi)$). Mass arises as closure residual: the photon has $\Delta = 0$ (perfect closure); the $W$ and $Z$ have $\Delta > 0$ (incomplete closure). The Weinberg angle is interpreted as a projection-induced rotation between torsional basis modes. The relation $m_W = m_Z \cos\theta_W$ is structurally preserved.

\textbf{Correspondence status:} Structural for the closure-residual reinterpretation of massless and massive modes; analogical for the mapping of gauge structure to torsional symmetry. The framework accommodates electroweak structure as a compatibility condition, but does not derive the gauge group.

%==================================================================
\section{Strong Interaction}\label{sec:strong}
%==================================================================

\textbf{Standard Model:} The strong interaction is governed by $\mathrm{SU}(3)$ colour gauge theory. Quarks carry colour charge; confinement prevents isolated colour-charged states from appearing as free particles. Hadrons are colour-neutral composites.

\textbf{Constraint framework (Paper~VIII~\citep{SmartVIII}):} Confinement is reinterpreted as a connectivity constraint: states with disconnected shell support are excluded from $\Scl$ by the gap penalty. The up quark on $\{2,4\}$ is disconnected and not closure-stable; the proton on $\{2,3,4\}$ is connected and closure-stable. Baryon-like composites involve three consecutive interior shells, providing a structural parallel to the threefold colour structure. $\mathrm{SU}(3)$ is not derived.

\textbf{Correspondence status:} Structural for confinement (connectivity = no free quarks). Analogical for colour (threefold shell structure $\leftrightarrow$ colour charge). $\mathrm{SU}(3)$ gauge dynamics are not reproduced.

%==================================================================
\section{Geometry and Gravity}\label{sec:gravity}
%==================================================================

\textbf{General relativity:} Spacetime is a pseudo-Riemannian manifold. Free particles follow geodesics. Curvature is sourced by mass-energy through the Einstein field equations.

\textbf{Constraint framework (Papers~IX--X~\citep{SmartIX, SmartX}):} The constraint manifold $\Mcl$ is equipped with an induced Riemannian metric on its smooth strata. Closure-compatible motion is represented in geodesic-like terms. Curvature arises from constraint interaction, not from a stress-energy source. Local flatness on smooth strata provides an analogue of locally inertial frames.

\textbf{Correspondence status:} Analogical throughout. The geometric form matches (curved manifold, geodesics, curvature), but the physical content differs: $\Mcl$ is a constraint manifold in state space, not a spacetime manifold. No Einstein equations, gravitational coupling constants, or cosmological solutions are derived.

%==================================================================
\section{Where the Framework Differs}\label{sec:differs}
%==================================================================

The constraint framework differs from the Standard Model in fundamental ways that must be stated explicitly:

\begin{itemize}
    \item \textbf{No Lagrangian.} The framework does not provide a Lagrangian density or action principle. The closure functional $\Fcal$ plays a structural role (characterising admissible states) but is not a dynamical action.
    \item \textbf{No field equations.} There are no equations of motion, no propagators, and no Feynman rules.
    \item \textbf{No renormalisation.} The framework operates at the level of constraint structure, not perturbative quantum field theory. Divergences, running couplings, and renormalisation group flow are not addressed.
    \item \textbf{No scattering amplitudes.} Cross-sections, decay rates, and transition amplitudes are not computable within the current framework.
    \item \textbf{No experimental predictions beyond mass correspondences.} The framework produces mass-level numerical correspondences (Paper~V) and structural compatibility with known features (confinement, electroweak relations), but does not yet generate experimentally distinguishable predictions.
    \item \textbf{No spacetime.} The constraint manifold $\Mcl$ is not identified with physical spacetime. The gravitational analogy (Paper~X) is structural, not physical.
\end{itemize}

\begin{remark}
These are not claimed as strengths of the framework. They are honest statements of what is missing. The framework in its current form addresses \textit{structural} questions (which states are admissible, what mass do they carry, how are they organised) but not \textit{dynamical} questions (how do they evolve, scatter, decay). A full physical theory would require both.
\end{remark}

%==================================================================
\section{Why This Might Matter}\label{sec:motivation}
%==================================================================

Despite these limitations, the constraint framework offers several features that may be of interest:

\begin{enumerate}
    \item \textbf{Parameter reduction.} The Standard Model has 19+ free parameters. The constraint framework replaces 13 of these (particle masses) with spectral-geometric correspondences from a single mathematical object (the 600-cell), using zero fitted continuous parameters.

    \item \textbf{Structural unification.} Mass, electroweak structure, confinement, and a gravitational analogy are all expressed within a single constraint-manifold formalism, rather than requiring separate mechanisms for each.

    \item \textbf{Alternative geometric starting point.} The framework begins with constraints and admissibility rather than fields and dynamics. This may provide a complementary perspective on why the Standard Model has the structure it does.

    \item \textbf{Falsifiability.} The mass correspondences are specific numerical predictions. Any deviation from experiment at improved measurement precision would constrain or falsify the framework.
\end{enumerate}

These are motivations, not proofs. The framework must ultimately be judged by whether it can be extended to make dynamical predictions that agree with experiment.

%==================================================================
\section{Open Questions}\label{sec:open}
%==================================================================

\begin{enumerate}
    \item \textbf{Why the 600-cell?} The privileged role of the 600-cell is supported by uniqueness results and eigenvalue correspondences, but not derived from a deeper principle.
    \item \textbf{Can $\mathrm{SU}(3)$ be derived?} The threefold composite structure is a geometric observation, not a gauge-group derivation.
    \item \textbf{Can dynamics be introduced?} The framework currently lacks evolution equations. Introducing a Lagrangian or Hamiltonian compatible with the constraint structure is a major open problem.
    \item \textbf{What is the relation to QFT?} The constraint framework operates at the level of state-space geometry. Its relationship to operator algebras, renormalisation, and perturbative QFT is not established.
    \item \textbf{Can curvature become gravity?} The gravitational analogy (Paper~X) is structural. Whether it can be developed into a physical theory of gravity remains open.
    \item \textbf{Can the framework predict new particles?} The mass correspondences are for known particles. Whether the framework predicts additional states in currently unexplored mass ranges is not yet determined.
\end{enumerate}

%==================================================================
\section{What Is and Is Not Claimed}\label{sec:claims}
%==================================================================

\subsection*{Claimed}

\begin{itemize}
    \item A systematic structural correspondence exists between the constraint framework and Standard Model concepts.
    \item Multiple physical sectors (mass, electroweak, strong, gravity) map consistently within this correspondence.
    \item The mass sector produces high-accuracy numerical correspondences with zero fitted parameters.
    \item The framework provides a unified geometric language for multiple physical phenomena.
\end{itemize}

\subsection*{Not claimed}

\begin{itemize}
    \item Derivation of the Standard Model.
    \item Physical equivalence between the constraint framework and the Standard Model.
    \item Derivation of general relativity.
    \item A complete physical theory (dynamics, scattering, renormalisation are absent).
    \item Experimental predictions beyond mass-level correspondences.
\end{itemize}

%==================================================================
\section{Conclusion}\label{sec:conclusion}
%==================================================================

The constraint-manifold framework developed in Papers~V--X admits a systematic structural correspondence with Standard Model concepts. Physical states correspond to the constraint manifold; gauge symmetry to projection invariance; mass to closure residual; confinement to connectivity constraints; and gravitational structure to manifold curvature. The correspondence is structural for the mass and strong sectors, analogical for the electroweak and gravitational sectors, and open for dynamics, renormalisation, and scattering.

This paper does not claim that the constraint framework derives or replaces the Standard Model. It provides a translation layer that makes the framework legible to physicists, identifies where the structural correspondence is strongest, and explicitly marks where it remains incomplete. The framework's value lies in its ability to organise multiple physical phenomena within a single constraint-geometric formalism using zero fitted continuous parameters. Whether this structural organisation reflects a deeper physical principle, or instead collapses under tighter dynamical and experimental tests, is the central question the programme must ultimately answer.

%==================================================================
\bibliographystyle{plainnat}
\bibliography{references}

\end{document}
